% !TEX root = ../metatime_monograph.tex
\chapter{Gauge Boson Masses}
\label{ch:gauge_bosons}

% CITATION_CONTEXT_V10_9_BEGIN
\paragraph{Established context.}
Non-Abelian gauge theory and electroweak unification originate with Yang--Mills
and Glashow--Weinberg--Salam \cite{YangMills1954,Glashow1961,Weinberg1967,Salam1968}.
Mass generation uses the established symmetry-breaking mechanism
\cite{EnglertBrout1964,Higgs1964,GHK1964}, and numerical boson properties are
benchmarked against PDG data \cite{PDG2024}.
% CITATION_CONTEXT_V10_9_END

\claimstatus{Frozen structural electroweak relation; precision comparison currently FAIL and requires a matched radiative/RG map.}

\section{Frozen SU(2) coupling relation}

The current structural relation is

\begin{equation}
g_0 = \frac{L_4}{L_3} + \frac{L_4}{L_5}
  = \frac{2}{7} + \frac{2}{5}
  = \frac{24}{35}
  = 0.68571.
\label{eq:g_coupling}
\end{equation}

The first term $L_4/L_3$ is the $\hat I$--$\hat M$ edge of the NOEMA
tetrahedron (colour--annihilation mixing), while $L_4/L_5$ is the
annihilation--intention mixing edge.  In the present monograph this quantity is
the frozen bare/structural electroweak coupling candidate $g_0$.

\section{Tree-level mass map}

The current tree-level map is

\[
M_W^{(0)} = \frac{g_0}{2}\,v.
\]

With $v = 244.89$~GeV (Chapter~\ref{ch:higgs_vev}),

\[
M_W^{(0)} = \frac{0.68571}{2}\times244.89
          = 83.96\ \gev.
\]

Using the weak-angle value from Chapter~\ref{ch:weinberg_angle},

\[
M_Z^{(0)} = \frac{M_W^{(0)}}{\cos\theta_W}
          = \frac{83.96}{0.87669}
          = 95.77\ \gev.
\]

\section{Frozen reference comparison}

For the reference values used in the current validation snapshot,

\[
M_W^{\rm ref}=80.377\ \gev,
\qquad
M_Z^{\rm ref}=91.188\ \gev,
\]

the fractional residuals are

\[
\frac{83.96-80.377}{80.377}=+4.46\%,
\]

and

\[
\frac{95.77-91.188}{91.188}=+5.02\%.
\]

The ratio residual is substantially smaller:

\[
\frac{M_Z^{(0)}}{M_W^{(0)}}
=\frac1{\cos\theta_W}
=1.1407,
\qquad
\frac{91.188}{80.377}=1.1346.
\]

This isolates the dominant discrepancy in the common electroweak scale/coupling
layer rather than in the ratio alone.

\section{PDG-2026 precision status}

The current $M_W^{(0)}$ and $M_Z^{(0)}$ values are retained as a frozen
structural/tree-level result.  Their several-percent residuals place this
version in the precision-failure class of the PDG-2026 validation ledger.

Precision promotion requires a single declared comparison scheme and scale for
$g$, $\sin^2\theta_W$, $v$, $M_W$ and $M_Z$, together with the corresponding
radiative or renormalization-group transport from the structural relation to
the observable pole quantities.

The required next map is therefore

\[
\boxed{
(g_0,\theta_W^{(0)},v_0)
\xrightarrow{\;\mathcal R_{EW}(\mu,\text{scheme})\;}
(g(\mu),\theta_W(\mu),v(\mu))
\longrightarrow
(M_W^{\rm pole},M_Z^{\rm pole}).
}
\]

The previously displayed multiplicative $g(1-\kappa)$ adjustment is removed
from the active repair path.  Future electroweak corrections enter through the
explicit matched map $\mathcal R_{EW}$ with frozen provenance before numerical
comparison.
